\scanpage{348}
\(\tfrac12\)-Algèbre différentielle \(C^*\) sur Topos \(\unsure{\mathscr X|X|\mathcal X}\)

Anneau de base \(k\) (commutatif)

résolution \(C^*\) de \(k\)
\margin{\unsure{ou, il suffit qu'il y ait une} augmentation}
% slanted
\diagram{\begin{tikzcd}[ampersand replacement=\&,column sep=small,row sep=small]
 \cdots\arrow[r,dash] \& C^0\arrow[r,dash] \& C^1\arrow[r] \& C^2\arrow[r] \&\cdots\\
 {} \& k\arrow[u] \& {} \& {} \&
\end{tikzcd}}

(la section \unsure{unité} de \(k\) correspond à une section \(T\) de \(C^0\)).

Accouplements \deleted{\unsure{sections}} (notés \(x,y\longmapsto x\circ y\))
\[
 \left\{\begin{aligned}
 Z^i\times\deleted{\unsure{Z}}C^j&\longrightarrow C^{i+j} &(a_{ij})\\
 C^i\times Z^j&\longrightarrow C^{i+j} &(b_{ij})
 \end{aligned}\right.\qquad i,j\geq0.
\]
satisfaisant les relations

\(1^\circ)\) \(a_{ij}\) et \(b_{ij}\) coïncident sur \(Z^i\times Z^j\) et y prennent valeurs dans \(Z^{i+j}\). (Donc la notation \deleted{\unsure{est}} \(x\circ y\) non ambigüe, elle a un sens si \(dx\) \uline{ou} \(dy\) est nul.)

\(2^\circ)\) Associativité \((x\circ y)\circ z=x\unsure{\circ|\,}(y\circ z)\), si deux des trois sections locales \(x,y,z\) (de \(C^i,C^j,C^h\)) sont des cocycles ; \(1\circ x=x\circ1=x\).

\(3^\circ)\)
\[
 \left\{\begin{aligned}
 x\circ y&=y\circ x\,(-1)^{\deg y\deg x}\\
 x\circ x&=0\quad\text{si \(x\) section locale de \(Z^i\), \uline{\(i\) impair}.}
 \end{aligned}\right.
\]
\inserted{Si \(x^i\) section locale de \(C^i\), \(z^j\) section locale \unsure{fermée}}
\deleted{\(4^\circ)\) Pour tout \(i,j\geq1\), existe \unsure{une} \(k\)-\unsure{bilin.} \(Z^j\), \unsure{on exige que} \(d\)}
\[
 \deleted{C^{i-1}\times Z^j\xrightarrow{b_{ij}} C^{i+j-1}}
 \quad \deleted{\text{\unsure{rendant le carré suivant commutatif}}}
\]
\diagram{\begin{tikzcd}[ampersand replacement=\&,column sep=large]
 B^{i}\times Z^j\arrow[r] \& Z^{i+j}\\
 C^{i-1}\times Z^j\arrow[u,"d\times\mathrm{id}"]\arrow[r,"b_{ij}"'] \&
 C^{i+j-1}\arrow[u,"d"']
\end{tikzcd}}
\[
 \boxed{d(x^i\circ z^j)=dx^i\circ z^j}
\]
\unsure{Alors}, \unsure{utilisant \(3^\circ\), on a} \(\unsure{d=(-1)^i z^i\circ dy^j}\).
\[
 \deleted{\unsure{d x^{(i)}|d(x^{(i)})}\times Z^j=d\bigl(b_{ij}(x^{i-1},z^j)\bigr)}
\]
\unsure{(i.e.} \(d(x^i\circ y^j)\) \unsure{est défini s'il s'agit de \(x^i\) ou \(y^j\) un cocycle}).

\deleted{(\uline{NB} \unsure{l'indétermination} des \(b_{ij}\) est dans \(\operatorname{Hom}_k(C^{i-1}\otimes_k Z^j,Z^{i+j-1})\) ; -- elles ne sont pas nulles ; on peut exiger que les \(\unsure{b_{ij}}\) soient \uline{données}, et pour des \unsure{raisons dessus} \(\ldots\))}
% archive sheet number: 349
